\section{Introduction}
Durant mes trois années d’études au sein du département informatique de l'INSA Rennes, je me suis particulièrement intéressé à la cybersécurité, tant dans ses aspects applicatifs qu'infrastructurels. J’ai choisi l’option sécurité lors du deuxième semestre de ma première année, ce qui m’a permis d’acquérir de solides connaissances dans ce domaine à travers des enseignements tels que : sécurité offensive, sécurité des réseaux, programmation sécurisée et cryptographie.

Pour valider ma dernière année du cycle d’ingénieur, j’ai opté pour un contrat de professionnalisation incluant un projet de fin d’études (PFE) tout au long de l’année. J’ai ainsi eu l’opportunité de réaliser cette alternance au sein de l’entreprise Orange Cyberdefense, en tant qu’ingénieur sécurité dans l’équipe Next Generation Firewall. Cette expérience m’a permis de développer davantage mes compétences et de renforcer mes connaissances dans le domaine qui me passionne.

Ce rapport présente les objectifs et les travaux réalisés durant cette période d’alternance. Dans un premier temps, je présenterai l’entreprise Orange Cyberdefense, suivie d’une description du contexte dans lequel s’est déroulé l'alternance. Je décrirai ensuite le (TODO). Par la suite, j’aborderai en détail les missions effectuées, les difficultés rencontrées ainsi que les solutions mises en œuvre. Enfin, je conclurai par une synthèse des principaux éléments abordés dans ce rapport.